% GLS-2026-001 rigour-without-infrastructure -- concept paper.
% Two-column arXiv-style preprint, themed after sources/CPRMH_arxiv_twocolumn_v12_core_collapse.pdf
% page 1 (grey small-caps running header, dark-blue title block, boxed meta line, shaded "In brief"
% box) per founder ruling BBL-2026-09-04-099 (arXiv two-column theme, polished). Built from
% templates/paper/arxiv-twocol/{main.tex,disclaimers.tex,refs.bib}; every mandatory template
% element (standpoint, tier legend, five-questions crosswalk, disclaimer catalogue, AI-assistance
% disclosure, claim matrix, Blackbox Note appendix, Human Mastery Gate, provenance) is kept.
% Genre: CONCEPTUAL, adopted-by-human override (BBL-2026-09-04-097; router's own mechanical output
% was systematic_review, not adopted -- see genre_decision.yaml).

\documentclass[10pt,twocolumn]{article}

\usepackage[margin=0.75in,columnsep=0.25in]{geometry}
% Language note: this paper is English only, no Thai edition, by founder
% ruling BBL-2026-09-04-103. Every Blackbox Note line quoted in the body is
% given as an English gloss (the AI assistant's translation, tier Dr, not
% the founder's own words) with a pointer to the verbatim Thai record --
% the Blackbox Log (DOI 10.5281/zenodo.22302518) or the cited note-file
% line -- rather than reproduced here in Thai.
\usepackage[T1]{fontenc}
\usepackage[utf8]{inputenc}
\usepackage{mathptmx}
\usepackage{microtype}
\usepackage{amsmath,amssymb,mathtools}
\usepackage{natbib}
\usepackage{booktabs}
\usepackage{tabularx}
\usepackage{array}
\usepackage{multirow}
\usepackage{enumitem}
\usepackage{graphicx}
\usepackage{xcolor}
\usepackage{url}
\usepackage[breaklinks=true]{hyperref}
% xurl: url.sty's default break points (mainly after '/') leave long
% hyphen-joined URL slugs (e.g. iseas.edu.sg's article-title-as-URL) as one
% unbreakable run under hyperref's \url; xurl allows breaks at any
% character so those DOI/URL bibliography entries no longer overflow the
% column width (confirmed: without it, one entry alone was 330pt overfull).
\usepackage{xurl}
\usepackage{cleveref}
\usepackage{etoolbox}
\usepackage{seqsplit}
% Long \texttt{...} citation-card filenames inside the reference list (e.g.
% cite-own-knowledge-stabilized-translation-h3-007.yaml) have no natural
% break point for plain \texttt -- a literal hyphen is not a break point in
% TeX's own paragraph builder for tt text -- so one entry alone was 45pt
% overfull. Scoped ONLY to thebibliography (natbib/plainnat's own generated
% .bbl markup, left untouched) so nothing outside the reference list changes.
\AtBeginEnvironment{thebibliography}{%
  \let\RWIOrigTexttt\texttt
  \renewcommand{\texttt}[1]{\seqsplit{#1}}%
}
\usepackage{titlesec}
\usepackage{caption}
\usepackage{float}
\usepackage{balance}
\usepackage{fancyhdr}
\usepackage{framed}

\definecolor{Ink}{RGB}{25,25,28}
\definecolor{Slate}{RGB}{82,88,96}
\definecolor{TitleBlue}{rgb}{0.10,0.20,0.45}
\definecolor{BriefBg}{RGB}{237,241,247}
\definecolor{MetaBg}{RGB}{247,247,249}
\definecolor{MetaRule}{RGB}{198,204,214}

\hypersetup{
  colorlinks=true, linkcolor=Ink, citecolor=TitleBlue, urlcolor=Slate,
  pdfauthor={Yaoharee Lahtee},
  pdftitle={Rigour Without Infrastructure: A Concept Paper (GLS-2026-001)}
}
\setlength{\parindent}{1em}
\setlength{\parskip}{0.2em}
\setlist{nosep,leftmargin=*}
\captionsetup{font=small,labelfont=bf}
\titleformat{\section}{\large\bfseries\scshape\color{TitleBlue}}{\thesection}{0.6em}{}
\titleformat{\subsection}{\normalsize\bfseries\color{TitleBlue}}{\thesubsection}{0.5em}{}

% ---- CPRMH-style running header/footer: small-caps grey preprint line ----
\pagestyle{fancy}
\fancyhf{}
\fancyhead[C]{\footnotesize\textsc{\color{Slate}Preprint \textbullet\ Concept paper \textbullet\ September 2026}}
\fancyfoot[C]{\footnotesize\color{Slate}\copyright\ 2026 The Author \textbullet\ CC BY 4.0 \textbullet\ glosa v0.2.0 arXiv two-column preprint \textbullet\ not yet peer reviewed $\mid$ \thepage}
\renewcommand{\headrulewidth}{0pt}
\renewcommand{\footrulewidth}{0pt}
\fancypagestyle{plain}{%
  \fancyhf{}
  \fancyhead[C]{\footnotesize\textsc{\color{Slate}Preprint \textbullet\ Concept paper \textbullet\ September 2026}}
  \fancyfoot[C]{\footnotesize\color{Slate}\copyright\ 2026 The Author \textbullet\ CC BY 4.0 \textbullet\ glosa v0.2.0 arXiv two-column preprint \textbullet\ not yet peer reviewed $\mid$ \thepage}
  \renewcommand{\headrulewidth}{0pt}
}

% ---- RWI tier-tag macros (readout_genesis + readout_universe tier ladder) ----
\newcommand{\ThCoqc}{\textsc{Th\_coqc}}
\newcommand{\FiniteDiag}{\textsc{finite\_diagnostic}}
\newcommand{\FitCal}{\textsc{fit\_calibrated}}
\newcommand{\DrTier}{\textsc{Dr}}
\newcommand{\DefTier}{\textsc{definition}}
\newcommand{\OpenTier}{\textsc{Open}}
\newcommand{\tier}[1]{\textsc{#1}}
\newcommand{\claimref}[1]{\textsuperscript{[#1]}}

% ---- Disclaimer-catalogue environment ----
\newlist{disclaimers}{description}{1}
\setlist[disclaimers]{style=nextline,leftmargin=1.4em,font=\bfseries}

% ---- Thin blockquote-style left rule for founder's-line English glosses
% (CPRMH-style quote rule; overrides framed.sty's default 3pt-wide bar) ----
\renewenvironment{leftbar}{%
  \def\FrameCommand{{\color{MetaRule}\vrule width 0.8pt}\hspace{8pt}}%
  \MakeFramed{\advance\hsize-1em \FrameRestore}}%
 {\endMakeFramed}

\begin{document}
\thispagestyle{plain}

% ---------------------------------------------------------------------------
% Full-width title / meta / "In brief" / abstract block (CPRMH page-1 theme)
% \twocolumn[...] keeps this front matter full-width at the top of page 1 and
% then resumes true two-column typesetting on the SAME page (a bare
% \onecolumn ... \twocolumn pair instead forces a page break before the
% column body, leaving page 1 half-empty -- avoid that).
% ---------------------------------------------------------------------------
\twocolumn[
\begin{center}
% Note: plain \\[dim] breaks inside a \twocolumn[...] optional argument
% (LaTeX kernel bug: \@icentercr's own optional-argument lookahead collides
% with \twocolumn's argument scanning -- confirmed by minimal reproduction);
% \par\vspace{dim} is the working equivalent for line breaks in this block.
{\LARGE\bfseries\color{TitleBlue} Rigour Without Infrastructure}\par\vspace{0.35em}
{\large\color{TitleBlue} Three Propositions on Claim-Card Discipline as a Substitute for Institutional Certification}\par\vspace{0.9em}
{\normalsize Yaoharee Lahtee}\par\vspace{0.15em}
{\small Open Civil Science Initiative, Thailand}\par\vspace{0.1em}
{\small ORCID \href{https://orcid.org/0009-0005-3861-0626}{0009-0005-3861-0626}}
\end{center}

\vspace{0.5em}
\begin{center}
\fcolorbox{MetaRule}{MetaBg}{%
  \begin{minipage}{0.94\linewidth}
  \small\centering
  \textbf{Article type} Concept paper (CONCEPTUAL genre, adopted-by-human override, \Cref{sec:genre})
  \textbullet\ \textbf{Method} Claim-card discipline + Literature Review System (LRS) self-application
  \textbullet\ \textbf{Licence} CC BY 4.0\par\vspace{0.25em}
  \textbf{Status} K0 working release, not peer reviewed
  \end{minipage}%
}
\end{center}

\vspace{0.6em}
\noindent\colorbox{BriefBg}{%
  \begin{minipage}{\linewidth}
  \small
  \textbf{In brief.} \textbf{Problem.} A person with no university, no lab, and no research team
  who now has AI assistance faces a new failure mode: text that reads as smart and academic while
  it stays unclear which part is evidence, which part is interpretation, and which part AI smoothed
  over past what the data supports (\Cref{sec:problem}). \textbf{Proposition.} Three candidate
  hypotheses (\Cref{sec:propositions}) --- a design mechanism (H1), a conceptual claim about where
  rigour lives (H2), and an empirical claim about cross-vendor AI checking (H3) --- propose that a
  claim-card discipline bound to Existence--Attribution--Disclosure and gated by an independence
  ladder could make a standalone scholar's claim reader-checkable without an institutional
  credential. \textbf{Method.} Each proposition carries a named falsifier and is placed against 48
  independently verified citation cards in a conversation-with-the-literature table
  (\Cref{sec:literature}), then the whole method is run on this paper's own construction as a
  self-application demonstration (\Cref{sec:demonstration}), including an independent review of the
  three claim cards themselves by three cross-vendor AI routes, which revised them
  (\Cref{sec:self-review}). \textbf{Status.} All three
  propositions are carried as untested claims at tier \DrTier{}, none above knowledge state K0; H3
  is explicitly untested, and the self-application run itself surfaced real gate failures, not a
  clean pass (\Cref{sec:demonstration}, \Cref{sec:limitations}).
  \end{minipage}%
}

\vspace{0.7em}
\begin{quote}
\small\noindent\textbf{Abstract.} People who live daily with a problem often know it before
researchers do, but turning that knowledge into a checkable claim has historically required a
university, a lab, or a research team. Generative AI removes that requirement technically while
introducing a new one epistemically: it becomes cheap to produce text that reads as rigorous
without it being clear which part is observation, which part is inference, and which part the AI
supplied. This paper states, as three falsifiable propositions and not as demonstrated results, a
candidate answer: a \emph{claim card} that separately records what was seen, what the data
separates, what AI filled in, what is assumed, and what independent check (if any) the claim has
survived --- bound to a formal Existence--Attribution--Disclosure norm and passed through a named
independence ladder (I0 self-check through I5 outside human review) before public release ---
could carry the checkable property institutions have historically supplied, without requiring the
institution itself (H1, design; H2, conceptual). A third, narrower proposition (H3, empirical)
holds, untested, that cross-vendor AI checking would measurably outperform same-vendor re-checking
at catching undisclosed AI-fill; no measurement of that difference has been made (n=0). Each
proposition is placed against a literature conversation built from 48 independently verified
citation cards, and the whole method is applied reflexively to this paper's own construction, which
surfaced real failures --- a validator warning, a manifest-detected search-mode gap, and a
documented citation-verification correction --- reported as findings, not smoothed into a clean
success story. The three claim cards carrying these propositions were themselves reviewed by three
cross-vendor AI routes and revised in response before this draft (\Cref{sec:self-review}). Every
claim here sits at tier \DrTier{} and knowledge state K0; none has passed an independent (I5) human
check.
\end{quote}
\vspace{0.7em}
]

\paragraph{Positioning.} We are not competing with knowledge-authority and make no priority claim;
we state what we propose, what we build on, and what would make us wrong.

% ===========================================================================
\section{Problem before observation}
\label{sec:problem}

``Observation'' is abstract; it is a verb, not a starting node. Founder's line, English gloss
below (Blackbox Log \texttt{BBL-2026-09-04-083}, concept DOI \texttt{10.5281/zenodo.22302518}):

% leftbar (framed.sty, trivlist-based) is used bare, not nested inside
% \begin{quote}: nesting two list-based environments here triggers a LaTeX
% "Something's wrong--perhaps a missing \item" error (confirmed by minimal
% reproduction) -- leftbar's own hsize-narrowing already gives the
% blockquote look on its own.
\begin{leftbar}\small
\textit{Founder's line (English gloss; verbatim Thai in the Blackbox Log, DOI
\texttt{10.5281/zenodo.22302518}, entry \texttt{BBL-2026-09-04-083}):}\par
\itshape\upshape --- ``the word `observation' is abstract, it is a verb; if we started from the \emph{problem}
instead, is there a process like this? We propose focusing on the problem first.''
\end{leftbar}

A problematic state exists in the world \emph{before} any knower acts on it --- before it is
looked at, named, or measured. The spine this paper's claims are built on (adopted this pass into
\texttt{design/FOUNDATION\_v0.5.md} \S2.1a, BBL-2026-09-04-084) therefore does not start with
``Reality $\rightarrow$'' as its first node --- the lens never grants direct access to reality,
only to what was read out ($R_A = O_A(W;\Pi_A) \neq W$) --- it starts one layer finer:

\begin{quote}\small
Problem state (unread) $\rightarrow$ first human readout (the Blackbox Note's raw line, verbatim,
untranslated) $\rightarrow$ human-language question $\rightarrow$ [lens-in] $\rightarrow$
hypothesis (lens-out, signed) $\rightarrow$ evidence $\rightarrow$ provisional knowledge (K0, ...)
\end{quote}

This paper's own problem state, before it was read out at all: a person who lives daily with a
problem, has no university appointment, no certified lab, and no standing research team, and now
has general-purpose AI assistance. The first human readout of that state (Thai, source of
truth), Blackbox Note line 1, English gloss below
(\texttt{blackbox/BLACKBOX\_NOTE\_glosa-paper\_2026-09-04.md\#L1}):

\begin{leftbar}\small
\textit{Founder's line (English gloss; verbatim Thai in
\texttt{blackbox/BLACKBOX\_NOTE\_glosa-paper\_2026-09-04.md\#L1}):}\par
\itshape\upshape --- ``how can a person with no university, no lab, and no research team produce
rigorous knowledge from their own site of practice?''
\end{leftbar}

The contaminated concept this question already flags: ``rigour'' ordinarily presupposes an
institutional mechanism (a lab, a department, a review committee) sitting behind a claim, rather
than naming a property a single claim can carry on its own, field by field.

% ===========================================================================
\section{Standpoint and ownership}
\label{sec:standpoint-and-ownership}
\label{sec:standpoint}

This paper is written from the standpoint of a social-enterprise practitioner and citizen (the
founder, Yaoharee Lahtee), not a substitute expert in computer science, philosophy of science,
library/information science, or research-integrity studies. \textbf{Disciplines not claimed:}
medicine, religious/fiqh authority, law, engineering, anthropology, political science, and formal
computer science beyond what is mechanically executed here. This is disclaimer
\textbf{D-STANDPOINT} / \textbf{D-NONEXPERT} in \Cref{sec:limitations}.

\subsection{Responsibility per arrow}
Every arrow on the spine above is named by who performed it. \textbf{Data $\rightarrow$
Inference} may be \texttt{human}, \texttt{ai}, or \texttt{joint}; the three project claim cards
behind this paper (\Cref{sec:claim-matrix}) record it as \texttt{joint} --- AI retrieved
literature and drafted candidate hypotheses, the founder framed the problem and question.
\textbf{Inference $\rightarrow$ Claim is always human.} AI is never an author. This is not a new
rule so much as this pass naming, per arrow, a distinction the repo already held at the
whole-artifact level (\texttt{AGENTS.md} gate rule 5).

\subsection{Ownership: problem, question, and the selection of the hypothesis stay human}
Founder's line, English gloss, first stated (\texttt{BBL-2026-09-04-086}) and then self-corrected
(\texttt{BBL-2026-09-04-088}, the correction shown here is the one adopted):

\begin{leftbar}\small
\textit{Founder's line (English gloss; verbatim Thai in the Blackbox Log, DOI
\texttt{10.5281/zenodo.22302518}, entry \texttt{BBL-2026-09-04-088}):}\par
\itshape\upshape --- ``if the problem is still ours, the question is still ours, and the \emph{selection} of
the hypothesis is still ours, and it leads to solving our problem, then use AI, use a cat, use a
bear, go ahead --- because that is the goal of knowledge: to solve a problem for someone.''
\end{leftbar}

Three acts on the spine are the human's: owning the \textbf{problem}, asking the
\textbf{question}, and \textbf{selecting} the hypothesis. Every other node --- retrieval, drafting,
analysis, checking, rendering --- may be performed by any tool. Concretely for this paper: AI
(the AI assistant) drafted candidate hypotheses H1, H2, and H3
(\texttt{hypotheses.md}); the founder's own act of selection, recorded in
\texttt{hypothesis\_selection.yaml} and confirmed by ruling \texttt{BBL-2026-09-04-095} (Thai
verbatim in \texttt{entries.jsonl}; gloss: ``bring all three together''), chose all three --- H1 and H3 are carried as untested propositions
with stated falsifiers, and H2 is advanced in the paper's own voice as its central conceptual
claim (\Cref{sec:genre}, \texttt{genre\_decision.yaml}).

% ===========================================================================
\section{The lens and the five questions}
\label{sec:lens}

\subsection{Claim-tier legend}
\label{sec:tiers}
Every claim in this paper carries exactly one tier tag; tiers are never collapsed or silently
upgraded (\Cref{tab:tiers}).

\begin{table}[t]
\centering
\scriptsize
\begin{tabularx}{\linewidth}{@{}l X@{}}
\toprule
\textbf{Tag} & \textbf{Meaning} \\
\midrule
\ThCoqc{} & Machine-checked, axiom-free formal result. \\
\FiniteDiag{} & Measured/computed on a finite, retained, reproducible input; not a limit or population claim. \\
\FitCal{} & A model calibrated against declared data; accuracy conditional on that fit. \\
\DrTier{} & Declared bridge / human narrative synthesis; a design or interpretive judgment, not a proof or measurement. \\
\DefTier{} & A stipulated definition used only where decision-bearing. \\
\OpenTier{} & Not established; a real open question, not smoothed over. \\
\bottomrule
\end{tabularx}
\caption{Claim-tier legend (readout\_genesis + readout\_universe ladder). Every propositional claim
in this paper (H1, H2, H3, \Cref{sec:propositions}) is \DrTier{}; the demonstration
(\Cref{sec:demonstration}) is \FiniteDiag{} where a command was actually executed.}
\label{tab:tiers}
\end{table}

\subsection{The readout-not-truth lens}
This paper's problem-to-hypothesis translation was performed under \emph{Readout Universe ---
Yaoharee Lahtee} (lens; DOI \href{https://doi.org/10.5281/zenodo.21529456}{10.5281/zenodo.21529456},
\href{https://doi.org/10.5281/zenodo.21665100}{10.5281/zenodo.21665100}; repositories
\url{https://github.com/morrocwi/readout_universe},
\url{https://github.com/morrocwi/readout_genesis}). Naming the lens credits its source and lets a
reader replace it; it does not make the lens, or any hypothesis drafted under it, true
(disclaimer \textbf{D-EXTERNAL-INPUT}). Everything read under this lens is a finite, retained
\emph{readout} of a source, never the source itself: a benchmark number, a citation's exact
passage, a claim card's own field, this paper's own prior draft --- each was produced by some
process, at some resolution, under some method, and is read as that readout, not as truth handed
down.

\subsection{Existence--Attribution--Disclosure and the five questions}
The Readout Condition is $E \wedge A \wedge D$: a claim satisfying Existence (some provenance path
for every load-bearing distinction), Attribution (identification of every non-source node), and
Disclosure (every AI-supplied distinction named, not silently absorbed into the claim) is what a
claim card is built to record, field by field, against the founder's own five questions
(\Cref{tab:five-q}).

\begin{table}[t]
\centering
\scriptsize
\begin{tabularx}{\linewidth}{@{}p{0.06\linewidth} X p{0.13\linewidth}@{}}
\toprule
\textbf{\#} & \textbf{Founder's question} & \textbf{Mechanically checkable?} \\
\midrule
Q1 & what was actually seen & presence only \\
Q2 & what the data separates & presence only \\
Q3 & what AI filled in & presence only \\
Q4 & what is assumed & presence only \\
Q5 & what independent evidence/objection & \textbf{the one genuinely enforceable check} \\
\bottomrule
\end{tabularx}
\caption{Founder's five questions $\Leftrightarrow$ E-A-D $\Leftrightarrow$ claim-card field group
(\texttt{design/FOUNDATION\_v0.5.md} \S3.1). A card that answers Q1/Q2/Q4 honestly but misattributes
Q3 --- crediting the source with a distinction AI actually supplied --- fails the Readout Condition
even though every field looks filled; \texttt{silent\_lift\_check} is the mechanized test for
exactly this. Thai originals: Blackbox Log \texttt{BBL-2026-09-04-083} ff.}
\label{tab:five-q}
\end{table}

Only Q5 --- whether an independent check was ever run, and at what independence class --- is
structurally enforceable by the kernel; Q1--Q4 are presence-checkable, never
correctness-checkable, because \emph{mechanical validity $\neq$ semantic validity}. This is why the
independence ladder (\Cref{sec:literature}, \Cref{sec:demonstration}) is the licensed route to
closing what the other four questions cannot mechanically close on their own.

\section{Genre and evidence standard}
\label{sec:genre}
This paper is genre-routed as \textbf{CONCEPTUAL} by a human, adopted-by-human override of the
router's own mechanical output (\texttt{genre\_decision.yaml}). The router (run with
\texttt{--search-log}) returned \texttt{systematic\_review} because a frozen \texttt{search\_log.yaml}
exists under each of \texttt{h1}/\texttt{h2}/\texttt{h3}; the founder's ruling
(\texttt{BBL-2026-09-04-097}, Thai verbatim in \texttt{entries.jsonl}; gloss: ``make this document an
academic article, a concept paper'') does not adopt that output --- a documented search log
supporting three candidate hypotheses is not the same fact as the genre of the claim the paper
itself advances, and the router has no way to see that distinction (\texttt{genre\_decision.yaml}
note). H2 (the conceptual, portable-property claim) is advanced in the paper's own voice; H1
(design) and H3 (empirical) are carried as propositions with stated falsifiers, neither run as a
study this pass. Evidence standard: every citation is a \texttt{status: VERIFIED} card
(\Cref{sec:literature}); every propositional claim stays \DrTier{} until an independent test says
otherwise.

% ===========================================================================
\section{Three propositions}
\label{sec:propositions}
Candidates were drafted by AI from the lens translation above; the founder selected among them
(\Cref{sec:standpoint-and-ownership}). All three are carried as propositions, not results.

\subsection*{H1 --- Design: the claim-card mechanism (\claimref{C1})}
\textbf{Statement.} If every claim a standalone scholar produces is recorded as a claim card
answering the five questions, tied to E-A-D, and passed through the independence ladder before any
public release, then a person with no institutional affiliation \emph{could, as a design
proposition,} co-produce knowledge with AI whose rigour a reader can check for themselves, without
an institutional credential serving as the certifying step. This is a proposed mechanism, not an
established one: this project's own single self-application (n=1, this paper's own production run,
\Cref{sec:demonstration}) is offered as one demonstration readout of the mechanism, not as a
general test across standalone scholars.

\textbf{Falsifier.} A claim card that passes every field the kernel can mechanically check is
nonetheless found by a genuine independent reviewer (I5) to have a false answer at Q2 or Q3 ---
showing the mechanical check does not track the epistemic property it claims to track. Not yet run:
this project has not secured an I5 (outside human, no stake) route, so this falsifier has not been
tested in either direction.

\textbf{What would count as evidence.} A corpus of claim cards run through the independence ladder
to at least I3, with a subset reaching genuine I5, compared against a matched corpus produced
without the discipline, scored on whether an independent reader can recover the
evidence/AI-fill boundary the maker recorded.

\subsection*{H2 --- Conceptual: rigour as a portable claim-level property (\claimref{C2})}
\textbf{Statement.} ``Rigour'' is not a property that institutional machinery (lab, department,
review committee) confers on a claim from the outside; it is proposed here, as a conceptual
contrast, to be \emph{at least in part} a property a single claim can carry in its own recorded
structure --- institutions are one historical mechanism that has produced this property, not shown
here to be the only possible one. This does not establish that institutions are non-load-bearing:
the cited evidence for it is mixed and mostly same-lineage (\Cref{tab:h2-lit}).

\textbf{Falsifier.} A claim card that answers all five questions honestly and completely, checked
at I3, is shown to still require an institutional credential --- not just an independent human
reviewer --- to be treated as rigorous by its actual intended readers, for at least one class of
reader/decision this project cares about.

\textbf{What would count as evidence.} Reader studies on whether a target audience (journal
editors, grant reviewers, a practitioner's own community) treats a well-formed claim card as
sufficient without also asking who backs it.

\subsection*{H3 --- Empirical: cross-vendor checking raises AI-fill detection (\claimref{C3}, untested)}
\textbf{Statement.} Hypothesis, not yet tested: when claim cards are checked by a different AI
vendor from the one that produced them (cross-vendor, independence level I3) instead of being
re-checked by the same vendor that produced the card (I0/I1), the rate at which undisclosed AI-fill
(a false Q3 field) or undisclosed assumptions (a false Q4 field) is caught \emph{would rise}
measurably compared to same-vendor re-checking. No seeded-error benchmark has been run by this
project (n=0 own trials); the single most directly on-point source found the effect
direction-dependent, not uniform (\Cref{tab:h3-lit}).

\textbf{Falsifier.} A controlled comparison (claim cards seeded with a known number of undisclosed
AI-fill/assumption instances, same-vendor vs.\ cross-vendor routes) shows no significant
difference, or shows cross-vendor performs worse.

\textbf{What would count as evidence.} A seeded-error benchmark with a pre-registered scoring
rubric; tiered \FitCal{} at best until run, \DrTier{} until then. \textbf{This paper does not run
that benchmark; H3 is stated, not tested.}

All three propositions were drafted under the same lens (\Cref{sec:lens}); none has an evidence
relation above independence class I3, and none of the underlying claim cards has yet been approved
by the founder for release (\texttt{independent\_check.status: PENDING} on all three --- checked by
cross-vendor routes R1--R3 and revised accordingly, not yet re-checked by a new route or
founder-approved, \Cref{sec:claim-matrix}).

% ===========================================================================
\section{Conversation with the literature}
\label{sec:literature}
Each table below is a \emph{conversation with the problem}, not a chronology or priority ranking
(\texttt{design/S14\_literature-review-system.md}): sources are placed by how they talk to the
proposition, never by who came first. Every row resolves to a citation card with
\texttt{status: VERIFIED} in the corresponding \texttt{litreview\_manifest.yaml}
(\texttt{records/lit/rigour-without-infrastructure/h\{1,2,3\}/}), meaning both
\texttt{metadata\_verified} (Crossref/OpenAlex/arXiv lookup) and \texttt{claim\_match\_verified}
(a decorrelated cross-vendor AI route, \texttt{route:cross-vendor-ai}, I3) hold. \textbf{Relation column}
follows the repo's comparison rule (same/different/cited, never a priority claim): \emph{same} =
the source's own finding is compatible with the proposition's direction; \emph{different} = the
source's own finding counts against it; \emph{cited} = the source bears on the topic without a
same/different verdict (orthogonal or undetermined in the underlying dialogue table); rows marked
\emph{own lineage, not independent} are the founder's own prior published pieces, collided against
world literature per instruction \texttt{BBL-2026-09-04-098} and never counted as independent
support.

\begin{table*}[t]
\centering
\scriptsize
\begin{tabularx}{\linewidth}{@{}p{0.145\linewidth} p{0.19\linewidth} p{0.11\linewidth} X@{}}
\toprule
\textbf{Citation card} & \textbf{Source} & \textbf{Relation} & \textbf{Bearing} \\
\midrule
\cite{cite-cstp-boundary-work-001} & Mayes 2022, Citizen Science and Scientific Authority & different & Moderate, scoped: credentialed-expert legitimacy-shaping work continues to matter in practice --- cross-vendor review narrowed this to bear only on credentialed legitimacy/acceptance pressure, not to be read as evidence the full claim-card + E-A-D + independence-ladder mechanism was tested and found wanting. \\
\cite{cite-disclosure-not-documentation-001} & Cwik 2026, Disclosure is not documentation & same & Weak, component-only: supports the H1 component that AI-assisted work needs structured documentation beyond disclosure --- cross-vendor review downgraded this from moderate to weak and scoped it to \texttt{SUPPORTS\_COMPONENT\_ONLY}, not evidence the combined mechanism was tested. \\
\cite{cite-llm-hallucination-survey-001} & LLM hallucination survey, arXiv:2510.06265 & cited & Shows AI cannot be trusted to flag its own errors --- exactly why H1 requires an external check rather than AI self-report. \\
\cite{cite-ocsdnet-honf-001} & OCSDNET/HONF, Open Hardware and Citizen Science (Indonesia, Thailand, Nepal) & same & Grounds H1's premise that non-institutional knowledge production is a live regional practice, not a hypothetical. \\
\cite{cite-own-blackbox-log-h1-006} & Blackbox Log (Lahtee, Zenodo) & own lineage, not independent & Documents the append-only raw-record practice H1 assumes is followed; no claim about rigour by itself. \\
\cite{cite-own-readout-condition-h1-002} & The Readout Condition (Lahtee) & own lineage, not independent & Supplies the E-A-D norm H1's claim card operationalizes; does not itself specify the card's fields. \\
\cite{cite-own-readout-genesis-h1-003} & Readout Genesis Standalone Synthesis (Lahtee) & own lineage, not independent & Supplies the independence-ladder vocabulary H1 tracks; does not itself define a claim-card schema. \\
\cite{cite-own-standalone-scholar-h1-005} & The Standalone Scholar (Lahtee) & own lineage, not independent & Names the dual-track conversion (AI-assisted formation $\to$ independent human friction) H1's card is scoped inside, not a replacement for. \\
\cite{cite-own-written-ai-h1-001} & Written by AI. Still True. (Lahtee) & own lineage, not independent & Names ``knower fetishism''; H1's card targets exactly the constitution-level property this paper argues authority derives from. \\
\cite{cite-own-zero-readout-h1-004} & What a Zero Readout Certifies (Lahtee) & own lineage, not independent & Illustrates machine-checked discipline; H1's own mechanical checks (schema validation) are markedly weaker. \\
\cite{cite-pgph-citizen-science-sea-001} & Kumar et al.\ 2025, Community perceptions of citizen science (South/SEA) & cited & Governance-inclusion framing, orthogonal to H1's rigour-checkability framing. \\
\cite{cite-rsos-citizen-epistemology-001} & Jaeger et al.\ 2023, An epistemology for democratic citizen science & same & Weak, general-context only: robust knowledge can arise through plural/citizen epistemology outside narrow institutional credentialing --- cross-vendor review confirmed this as \texttt{SUPPORTS\_GENERAL\_CONTEXT} background support, not upgraded to a test of H1's mechanism. \\
\cite{cite-rsos-preprint-credibility-001} & Soderberg, Errington \& Nosek 2020, Credibility of preprints & same & Researchers rate content/process cues over prestige when prestige is absent --- supports the premise, not H1's specific design. \\
\cite{cite-scholarly-kitchen-credentialism-001} & Scholarly Kitchen 2026, The Guild and the Gap & cited & A real rejection episode on prior-publication grounds; whether a claim card would change the outcome is untested. \\
\cite{cite-tcij-freelance-th-001} & TCIJ (TH) 2020, freelance-academic labour piece & cited & Adjacent labour/credit question, thinner on H1's actual rigour-verification topic than its search match suggested. \\
\bottomrule
\end{tabularx}
\caption{H1 (design proposition) against 15 verified citation cards, 6 of them own-lineage.
Diversity audit (\texttt{h1/litreview\_manifest.yaml}): 1 Thai / 5 English external sources
(own-lineage excluded), stance counts \emph{agrees}:4 \emph{orthogonal}:2 \emph{undetermined}:9 by
the manifest's own collapsing convention (\Cref{sec:demonstration} notes a discrepancy between this
convention and the citation cards' own explicit notes). \texttt{concentration\_flags: []};
\texttt{zero\_disagree\_flag: false}.}
\label{tab:h1-lit}
\end{table*}

\begin{table*}[t]
\centering
\scriptsize
\begin{tabularx}{\linewidth}{@{}p{0.16\linewidth} p{0.19\linewidth} p{0.1\linewidth} X@{}}
\toprule
\textbf{Citation card} & \textbf{Source} & \textbf{Relation} & \textbf{Bearing} \\
\midrule
\cite{cite-rigour-without-infrastructure-001} & Wikipedia, Open-notebook science & cited & Treats openness as a property of record-keeping practice, close to but coarser than H2's own five-field structure. \\
\cite{cite-rigour-without-infrastructure-002} & Rasmussen 2019, Confronting Research Misconduct in Citizen Science & cited & Calls for institution-like misconduct-detection machinery --- cuts against H2's claim that the property needs no such mechanism. \\
\cite{cite-rigour-without-infrastructure-003} & Freitag, Meyer \& Whiteman 2016, credibility-building strategies & same & Credibility built bottom-up through concrete recorded practices, not top-down certification. \\
\cite{cite-rigour-without-infrastructure-004} & Scholarly Kitchen 2026, The Guild and the Gap (+ comments) & different & Close to H2's own falsifier: a named commenter states affiliation/publication history earns \emph{more} scrutiny than any single record could substitute for. \\
\cite{cite-rigour-without-infrastructure-005} & Dutta et al.\ 2021, Decolonizing Open Science & cited & Structural/platform-governance framing, adjacent to but different from H2's single-claim scope. \\
\cite{cite-rigour-without-infrastructure-006} & Okafor et al.\ 2022, Institutionalizing Open Science in Africa & cited & Aims for \emph{more} institutional machinery as the path forward --- the opposite emphasis from H2, at a different level. \\
\cite{cite-rigour-without-infrastructure-007} & Kumar et al.\ 2025, community perceptions of citizen science & cited & Access to institutional audiences, not record quality, is the bottleneck this source finds --- adjacent friction H2 does not resolve. \\
\cite{cite-rigour-without-infrastructure-008} & Srirot \& Sohsomboon 2025 (TH), qualitative Trustworthiness framework & same & Record-level techniques (audit trail, triangulation) let a reader judge rigour without institutional standing --- the closest structural match to H2 found. \\
\cite{cite-rigour-without-infrastructure-009} & THE STANDARD (TH) 2025, Thai Open Science policy piece & cited & Institution-led opening of an existing record --- the institution still decides what counts as legitimate, an access question not a rigour-provenance one. \\
\cite{cite-own-blackbox-log-h2-006} & Blackbox Log (Lahtee) & own lineage, not independent & No claim about whether rigour is portable at claim level. \\
\cite{cite-own-civilization-of-knowledge-h2-007} & The Civilization of Knowledge (Lahtee) & own lineage, not independent & Locates the stabilizing step historically in institutions --- would ask whether a claim card is a new non-institutional stabilizer or a mistaken substitute. \\
\cite{cite-own-readout-condition-h2-002} & The Readout Condition (Lahtee) & own lineage, not independent & Weak, indirect: supports source-relative warrant architecture, not explicitly ``rigour'' or institutional non-conferment --- cross-vendor review downgraded this from moderate/direct to weak/indirect support; own lineage, not independent confirmation. \\
\cite{cite-own-readout-genesis-h2-003} & Readout Genesis Standalone Synthesis (Lahtee) & own lineage, not independent & Institutions are one belief-scale among several, not a precondition for knowledge status. \\
\cite{cite-own-standalone-scholar-h2-005} & The Standalone Scholar (Lahtee) & own lineage, not independent & Independent human friction is one specific historical route to stress-tested knowledge, not necessarily the only one --- would ask if H2 still needs some friction step. \\
\cite{cite-own-written-ai-h2-001} & Written by AI. Still True. (Lahtee) & own lineage, not independent & Institutional authority is derivative, not a source, of epistemic value --- the same possession/constitution distinction H2 draws. \\
\cite{cite-own-zero-readout-h2-004} & What a Zero Readout Certifies (Lahtee) & own lineage, not independent & Illustrates zero-institutional-apparatus rigour in a formal setting --- weak indirect support only. \\
\bottomrule
\end{tabularx}
\caption{H2 (conceptual proposition, this paper's own advanced claim) against 16 verified citation
cards, 7 own-lineage. Diversity audit (\texttt{h2/litreview\_manifest.yaml}): 3 Thai / 13 English,
stance \emph{agrees}:6 \emph{disagrees}:2 \emph{orthogonal}:7 \emph{undetermined}:1.
\texttt{concentration\_flags: []}; \texttt{zero\_disagree\_flag: false} --- this is the one
proposition of the three whose manifest itself records a nonzero disagree count.}
\label{tab:h2-lit}
\end{table*}

\begin{table*}[t]
\centering
\scriptsize
\begin{tabularx}{\linewidth}{@{}p{0.155\linewidth} p{0.19\linewidth} p{0.1\linewidth} X@{}}
\toprule
\textbf{Citation card} & \textbf{Source} & \textbf{Relation} & \textbf{Bearing} \\
\midrule
\cite{cite-selfpreference-judge-001} & Wataoka, Takahashi \& Ri 2024/25, Self-Preference Bias in LLM-as-a-Judge & cited & Weak, context-only: finds familiarity/self-preference bias, a mechanism relevant to why same-vendor re-checking may under-detect, but does not itself run a cross-vendor-vs-same-vendor comparison --- cross-vendor review downgraded this from moderate to weak, scope narrowed to \texttt{CONTEXT\_ONLY}/\texttt{SUPPORTS\_GENERAL\_MECHANISM}. \\
\cite{cite-crossmodel-codereview-002} & Xiang et al.\ 2026, Cross-Model LLM Code Review, arXiv:2607.21656 & different & The one source that actually runs H3's own comparison: the effect is asymmetric and direction-dependent (one direction helps a lot, the reverse hurts), and same-vendor re-review also improved one condition. \\
\cite{cite-toohconsistent-003} & Tan et al.\ 2025, Too Consistent to Detect (EMNLP) & same & Weak, context-only: same-model repeated sampling cannot catch ``self-consistent'' errors, relevant to but not a direct test of H3's cross-vendor mechanism --- cross-vendor review downgraded this from moderate to weak, scope narrowed to \texttt{CONTEXT\_ONLY}/\texttt{SUPPORTS\_GENERAL\_MECHANISM}. \\
\cite{cite-multiagent-debate-004} & Du et al.\ 2023, Multiagent Debate & cited & Multi-instance debate improves factuality without varying vendor identity --- the active ingredient could be plurality, not cross-vendor difference. \\
\cite{cite-selfcorrect-reasoning-005} & Huang et al.\ 2023/24, LLMs Cannot Self-Correct Reasoning Yet & cited & Tests zero-external-signal correction, a weaker condition than a same-vendor re-check with a fresh adversarial prompt --- an imperfect analogy, flagged as such. \\
\cite{cite-legalai-india-006} & John, Singh \& Panachakel 2026, Can Legal AI Know When It Is Wrong? & cited & Three vendors' error rates differ nearly 5x on identical inputs, consistent with H3's premise, but the paper's own lever for improved catching is human experience, not cross-vendor AI review. \\
\cite{cite-bangladesh-readiness-007} & Sultana et al.\ 2026, Bangladesh AI Readiness & cited & Names a capacity defeater of H3's practical reach (infrastructure to run any independent review at all), not of H3's empirical claim. \\
\cite{cite-iseas-sea-governance-008} & Fong 2026, AI Governance Policies in Southeast Asia & cited & Institutional-capacity constraint on deploying cross-vendor review as policy, orthogonal to H3's mechanism claim. \\
\cite{cite-cofact-thai-detector-009} & Cofact Thailand 2025, AI-image-detection reliability & cited & Adjacent domain (detecting whether content is AI-made, not whether AI-authored claims are honest) --- genuinely different problem. \\
\cite{cite-mindstudio-crossvendor-010} & MindStudio blog 2026, Cross-Vendor AI Agent Review & same & Weak, context-only: states H3's own mechanism in practitioner language, but with no measurement, sample, or method of its own --- cross-vendor review scoped this \texttt{CONTEXT\_ONLY}/\texttt{SUPPORTS\_GENERAL\_MECHANISM}, flagged as concentration risk, not equivalent-weight evidence. \\
\cite{cite-own-blackbox-log-h3-006} & Blackbox Log (Lahtee) & own lineage, not independent & No claim about detection rates. \\
\cite{cite-own-knowledge-stabilized-translation-h3-007} & Knowledge as Stabilized Translation (Lahtee) & own lineage, not independent & Stability under variation is the marker of knowing well --- conceptually the same move H3 makes; does not specify vendor identity as the variation dimension. \\
\cite{cite-own-readout-condition-h3-002} & The Readout Condition (Lahtee) & own lineage, not independent & Disclosure is structural, not optional; does not specify the auditor must be a different AI vendor. \\
\cite{cite-own-readout-genesis-h3-003} & Readout Genesis Standalone Synthesis (Lahtee) & own lineage, not independent & Does not separate same-vendor from cross-vendor checking as a variable. \\
\cite{cite-own-standalone-scholar-h3-005} & The Standalone Scholar (Lahtee) & own lineage, not independent & Independent validation in this paper's stronger sense requires human friction --- cross-vendor AI (I3) is a weaker approximation. \\
\cite{cite-own-written-ai-h3-001} & Written by AI. Still True. (Lahtee) & own lineage, not independent & Reinforces that a checker's identity/route matters to what its check can license --- no detection-rate comparison. \\
\cite{cite-own-zero-readout-h3-004} & What a Zero Readout Certifies (Lahtee) & own lineage, not independent & Machine-checked verification, where available, is a strictly stronger independence route than cross-vendor AI review. \\
\bottomrule
\end{tabularx}
\caption{H3 (empirical proposition, explicitly untested) against 17 verified citation cards, 7
own-lineage. Diversity audit (\texttt{h3/litreview\_manifest.yaml}): 1 Thai / 6 English,
stance \emph{agrees}:3 \emph{orthogonal}:4 \emph{undetermined}:10 by the manifest's own convention.
\texttt{concentration\_flags: []}; \texttt{zero\_disagree\_flag: false}.}
\label{tab:h3-lit}
\end{table*}

Across all three tables, 20 of 48 rows (42\%) are the founder's own prior published pieces,
collided against world literature per instruction \texttt{BBL-2026-09-04-098} rather than left
silent; every own-lineage row is marked \emph{own lineage, not independent} and none is counted
toward any of the three propositions' evidential support.

% ===========================================================================
\section{The artifact applied to itself}
\label{sec:demonstration}
This section reports what the 2026-09-04 self-application run --- glosa run on glosa's own paper,
per instruction \texttt{BBL-2026-09-04-091} --- actually recorded, as a readout of that run, not as
a claim that the method passed cleanly.

\subsection*{Gate results actually executed}
Running \texttt{./cli/glosa claim validate} directly against the three project claim cards
(\texttt{GLOSA-CC-20260904-0201/0202/0203}) this session returns \texttt{verdict:
PASS\_WITH\_LIMITS} for all three, with zero schema errors, but with one identical warning on every
card: \emph{``D-LENS-UNCITED not checked -- validate\_claim\_card was not given citation\_cards
(lens\_translation.lens\_ref is set); pass citation\_cards=[\ldots] to close this gap.''} The
validator is honest about its own coverage limit here: a schema-level PASS is not, by itself, a
claim that the lens-citation gate closed --- it closed only once this session additionally
cross-checked the lens DOIs against the citation-card apparatus described in \Cref{sec:lens}.

\subsection*{Composite quotes caught by the cross-vendor route}
The single strongest gate finding this run: an I3 (decorrelated cross-vendor) check on an earlier
citation pass caught \emph{composite quotes} --- citations assembled from memory rather than a
locator taken from the opened source. The founder's ruling in direct response, English gloss
(\texttt{BBL-2026-09-04-100}):
\begin{leftbar}\small
\textit{Founder's line (English gloss; verbatim Thai in the Blackbox Log, DOI
\texttt{10.5281/zenodo.22302518}, entry \texttt{BBL-2026-09-04-100}):}\par
\itshape\upshape --- ``should we make it a hard rule: if a source is used it must carry a page and line locator plus
a full link, never written from memory --- that reduces the problem at the source, before the
check.''
\end{leftbar}
\noindent Adopted immediately (\texttt{BBL-2026-09-04-101}, Thai verbatim in \texttt{entries.jsonl}; gloss: ``update it into the
system''), this became the rule under which every one of the 48 citation cards behind
\Cref{sec:literature} was subsequently built: each carries \texttt{page\_or\_locator} and
\texttt{fetched\_from\_url}, and each \texttt{exact\_passage} is a source-first quotation, not a
recalled paraphrase. This is a real, gate-caught failure that changed the method's own rule
(rule 17, \texttt{design/FOUNDATION\_v0.5.md} \S7.8), reported here because it happened, not
because it flatters the method.

\subsection*{A manifest-detected search-mode gap}
The genre router's own mechanical output (\Cref{sec:genre}) returned \texttt{systematic\_review}
because \texttt{context.has\_search\_log} was true --- but every one of \texttt{h1}/\texttt{h2}/\texttt{h3}'s
own \texttt{litreview\_manifest.yaml} records \texttt{search\_mode} as \texttt{TARGETED\_SEARCH} /
\texttt{SCOPING\_SEARCH}, never \texttt{SYSTEMATIC\_REVIEW} --- the manifest, read directly, is what
let the founder catch that the router's mechanical genre suggestion did not match the actual
evidentiary standard the literature runs had met, and overrule it (\Cref{sec:genre}). This is the
manifest doing exactly the job \texttt{design/FOUNDATION\_v0.5.md} \S7.8's \textbf{FC-S8-1} hard
block exists for: calling anything less than a full SR protocol a ``systematic review'' is blocked,
here by the human reading the manifest against the router's own claim.

\subsection*{A stance-counting discrepancy this session found and did not paper over}
Building \Cref{tab:h1-lit,tab:h3-lit} directly from each citation card's own \texttt{notes} field
(``Dialogue-table stance: NO (disagrees with / challenges H1)'' on \texttt{cite-cstp-boundary-work-001})
surfaced a discrepancy with the corresponding manifest's own \texttt{diversity\_audit.counts.by\_stance}:
the manifests for H1 and H3 both record \texttt{disagrees: 0}, while this session confirmed
directly (grep against the raw dialogue-table rows, \FiniteDiag{}) that at least one row per
hypothesis (Mayes 2022 against H1; Xiang et al.\ 2026 against H3) is an explicit, single-value
challenge row that a card's own notes field independently labels a disagreement. The manifest's
aggregate counting convention appears to collapse a dual-value ``NO/NO'' dialogue-table cell into
\texttt{undetermined} rather than \texttt{disagrees}, undercounting real challenges in its own
summary statistic by at least one row per hypothesis where this occurs. This paper's own tables
(\Cref{tab:h1-lit,tab:h3-lit}) use the citation card's explicit \texttt{notes} field, not the
manifest's aggregate count, and this discrepancy is disclosed rather than silently reconciled ---
a small instance of exactly the ``mechanical check does not track the epistemic property it claims
to track'' failure mode H1's own falsifier names, found here in the tool that mechanically
summarizes H1's own supporting literature.

\subsection*{Path leaks caught by CI, not by this run}
The repository's public-facing leak gate (\texttt{scripts/check\_leak.sh}, wired into
\texttt{.github/workflows/ci.yml}'s \texttt{leak-scan} job) is a repo-level, not a paper-level,
control: it did not need to run specifically against this paper's files this session (none of this
paper's content matched a denylisted pattern), but it is the mechanical backstop this paper's own
public-facing disclosures rely on, and its presence in CI --- rather than only as a manual
pre-publish step --- is itself part of what this paper's demonstration reports: a leak gate that
runs on every push, not only when someone remembers to ask for it.

\subsection*{Independent review of this paper's own claims}
\label{sec:self-review}
The three project claim cards behind this paper (\Cref{sec:claim-matrix}) were themselves put
through the independence ladder before this draft was finalized: three cross-vendor AI review
routes (R1 Falsifier, R2 HostileReviewer, R3 SourceAuditor; independence class I3), producing nine
review reports in total (three routes $\times$ three cards,
\texttt{projects/GLS-2026-001\_rigour-without-infrastructure/reviews/*.review\_report.json}), read
each card independently of the session that drafted it. All nine converged on one defect class:
statements written as results when the evidence behind them was one self-application run plus
same-lineage or adjacent-domain citations, not a direct test of the mechanism stated. Concretely:
H1's ``can co-produce knowledge'' was rewritten to ``could, as a design proposition,'' with the
paper's own n=1 self-application named explicitly as one demonstration readout, not a general test;
H2's categorical ``rigour is not conferred by institutions'' was rewritten to a hedged conceptual
proposition (``proposed here $\ldots$ at least in part''), with every same-lineage supporting
citation now carrying an explicit not-independent note; H3's asserted ``rises measurably'' was
rewritten to an untested hypothesis, ``would rise $\ldots$ no measurement has been made, n=0.''
Every individual finding from all three routes is logged as a resolved dissent record on its
originating claim card (\texttt{five\_questions.tested.dissent\_records}), not silently absorbed
into a cleaner-looking rewrite. As of this draft, each card's \texttt{independent\_check.status} is
\textbf{PENDING}: checked by cross-vendor routes and revised accordingly, but not yet re-checked by
a further independent route, and not yet approved by the founder for release
(\Cref{sec:claim-matrix}). This section is offered as a working example of the independence ladder
catching over-scope in this paper's own output --- one instance, at I3, on one paper --- not as
proof that the ladder works in general; that broader claim would itself need the same discipline
(H1's own falsifier, \Cref{sec:propositions}) before it could be asserted.

% ===========================================================================
\section{Limitations and disclaimers}
\label{sec:limitations}
\label{sec:disclaimers}
% disclaimers.tex — GLS-2026-001 rigour-without-infrastructure concept paper.
% Every id below is either (a) one of the ids emitted by
% `./cli/glosa claim disclaimers` against the three project claim cards
% (GLOSA-CC-20260904-0201/0202/0203, `disclaimers_emitted`, task B1) or (b) a
% template-mandatory id filled honestly for this paper. Ids are the semantic
% ids of the single disclaimer catalogue; no numeric D1-D12 scheme.

\begin{disclaimers}
  \item[D-STANDPOINT] See \Cref{sec:standpoint}. This paper is written from a
    social-enterprise-practitioner/citizen standpoint, not disciplinary
    authority; a reader who skips to this section should read every claim
    above as offered from that position and access level, never from
    institutional credential.
  \item[D-NONEXPERT] The author is not a substitute for a credentialed expert
    in citizen-science studies, philosophy of science, library/information
    science, or empirical AI-evaluation methodology --- the disciplines this
    paper draws on without claiming standing in. Medicine, religious/fiqh
    authority, law, engineering, anthropology, political science, and formal
    computer science beyond what is mechanically executed here are explicitly
    not claimed.
  \item[D-SCOPE] Not applicable in the empirical-population sense (this is a
    \textsc{CONCEPTUAL} genre paper, \Cref{sec:genre}). Where a source's own
    sample size or population matters to a comparison row
    (\Cref{tab:h1-lit,tab:h2-lit,tab:h3-lit}), it is stated inline in that
    row rather than in one blanket population statement.
  \item[D-NONCLAIM] This paper does \emph{not} claim: that H1's mechanism has
    been run as a study; that H2's conceptual-portability claim has been
    tested against real readers; that H3's seeded-error benchmark has been
    built or executed; that the 48 citation cards behind
    \Cref{tab:h1-lit,tab:h2-lit,tab:h3-lit} constitute a systematic or
    exhaustive review (\texttt{search\_mode: TARGETED\_SEARCH} /
    \texttt{SCOPING\_SEARCH}, never \texttt{SYSTEMATIC\_REVIEW}); that any
    finding here is a medical, legal, or religious determination; or that
    cross-vendor AI checking (I3) is equivalent to, or a substitute for, an
    I5 human review.
  \item[D-AIFILL] See \Cref{sec:ai-disclosure}. AI (the AI assistant, "the AI assistant" seat)
    drafted candidate hypotheses H1/H2/H3, the literature-review runs, the
    claim cards' \texttt{current\_evidence}, \texttt{retrieved\_tool\_evidence},
    \texttt{retained\_record\_route}, \texttt{model\_calibration\_assumption},
    \texttt{prompt\_system\_constraint}, and \texttt{decision\_policy} fields,
    and this paper's own prose and \LaTeX{} build. None of this is treated as
    evidence on its own say-so (\Cref{sec:ai-disclosure}).
  \item[D-TIER] The single highest tier claimed anywhere in this paper is
    \DrTier{} (declared bridge / human narrative synthesis over a specified
    mechanism), with one \FiniteDiag{} exception: the mechanically executed
    checks reported in \Cref{sec:demonstration} (schema validation, disclaimer
    emission, manifest gate reads) and this paper's own \LaTeX{} compile. It
    cannot go higher because none of H1/H2/H3 has been run as a study, and no
    claim card behind this paper has passed a class-5 (I5) independent human
    check.
  \item[D-INDEPENDENCE] Every citation card behind
    \Cref{tab:h1-lit,tab:h2-lit,tab:h3-lit} was checked at independence class
    $I3$ (decorrelated cross-vendor AI route, \texttt{claim\_match\_verified\_by:
    route:cross-vendor-ai}) for \texttt{claim\_match}, and mechanically ($I3$-adjacent,
    Crossref/OpenAlex lookup) for \texttt{metadata\_verified}. The three
    project claim cards themselves (\Cref{sec:claim-matrix}) sit at
    \texttt{independent\_check.status: NONE} --- drafted and self-checked by
    the authoring AI session, never yet reviewed by a route independent of
    that session (this is $I0$ for the cards' own content, not to be conflated
    with the $I3$ citation-level checks that feed them). No external or
    institutional validation is sought or treated as available as a lever
    that would make any claim here more true; an outside reader's input, if
    it comes, is optional-level evidence at whatever independence class it
    actually reaches, never a certification gate.
  \item[D-DVP-NOT-K2] A decorrelated-vendor AI review pass (DVP) is at most an
    $I3$ route. \textbf{DVP $\neq$ K2}: no combination of AI-only checks,
    however many vendors, raises this paper's claims to K2. This paper's
    K-state is \texttt{K0} (\Cref{sec:kstate}); no claim in it is asserted
    above K0/\DrTier{}.
  \item[D-LEGAL-NEQ-EPISTEMIC] None identified --- no license, permit, or
    regulatory approval is referenced by this paper's claims. The general
    rule still holds: were one ever cited, it would make an action lawful,
    never raise epistemic warrant or certify a claim as true.
  \item[D-NOT-DIAGNOSTIC] Not applicable --- this paper makes no health,
    clinical, or treatment recommendation of any kind.
  \item[D-TRANSLATION] Not applicable: per founder ruling
    (BBL-2026-09-04-099, \Cref{sec:blackbox}), this concept paper is
    published in English only, with no Thai edition. Every founder line
    quoted in \Cref{sec:blackbox} nonetheless stays Thai, verbatim,
    untranslated in place, with an English gloss alongside it --- the
    English-only rule governs the paper's own prose, not the Blackbox Note's
    raw lines. Within the underlying claim cards, the Thai statement fields
    themselves carry \texttt{translation\_status: reviewed}.
  \item[D-DISSENT-PRESERVED] None recorded to date for H1/H2/H3 specifically.
    The literature-review dialogue tables (\Cref{tab:h1-lit,tab:h2-lit})
    already carry the closest real dissent this project has found against its
    own hypotheses (the Mayes 2022 boundary-work finding against H1
    \citep{cite-cstp-boundary-work-001}; the Scholarly Kitchen credentialism
    case against H2 \citep{cite-rigour-without-infrastructure-004}); these are
    stated as open challenges in \Cref{sec:literature}, not resolved or
    smoothed over.
  \item[D-COMPARISON] This paper does not claim priority or precedence.
    \Cref{sec:literature} states every source relation as \emph{agrees},
    \emph{disagrees}, or \emph{orthogonal/undetermined} --- never as a
    priority contest, never as ``no other work does X.''
  \item[D-CITATION-UNVERIFIED] None identified for the 48 citation cards cited
    in \Cref{tab:h1-lit,tab:h2-lit,tab:h3-lit} --- every one carries
    \texttt{status: VERIFIED} (\texttt{metadata\_verified: true},
    \texttt{claim\_match\_verified: true}) in its litreview manifest as of
    2026-09-04; none is cited from memory (\Cref{sec:demonstration},
    BBL-2026-09-04-100/101). Where a source's own full text was not fetched
    (e.g.\ \texttt{fetch\_status: PAYWALLED\_ABSTRACT\_ONLY}), the citation
    card's \texttt{scope} field names the resulting evidentiary weakness
    (\texttt{PARAPHRASE} rather than \texttt{DIRECT\_QUOTATION}) inline in
    the relevant table row rather than silently upgrading it.
  \item[D-BLACKBOX-NOTE] This paper carries the mandatory Blackbox Note
    appendix (\Cref{sec:blackbox}). Note id: \texttt{BB-2026-09-04-01}
    (\texttt{projects/GLS-2026-001\_rigour-without-infrastructure/blackbox\_note.yaml});
    curated public source:
    \texttt{blackbox/BLACKBOX\_NOTE\_glosa-paper\_2026-09-04.md}; public log
    entries \texttt{BBL-2026-09-04-083, -086, -088, -089, -091, -095, -097,
    -098, -099, -100, -101} (\texttt{blackbox/log/entries.jsonl}, concept DOI
    \texttt{10.5281/zenodo.22302518}).
  \item[D-K-STATE] Overall K-state: \texttt{K0} (internal working note). Every
    load-bearing claim in this paper --- H1, H2, H3, and every finding in
    \Cref{sec:demonstration} --- shares this same K0 state; there is no
    per-claim exception. This matches the three underlying project claim
    cards' own \texttt{k\_state: K0} (\Cref{sec:claim-matrix}).
  \item[D-AUTHORSHIP] See \Cref{sec:ai-disclosure}. The problem, the research
    question, and the selection among H1/H2/H3 are the founder's own act
    (\Cref{sec:standpoint-and-ownership}); AI drafted every candidate,
    citation search, and this paper's prose, disclosed per claim
    (\texttt{produced\_by}, \Cref{tab:claim-matrix}), and never substitutes
    for the author's judgment or responsibility
    (\textbf{AIContribution $\neq$ EpistemicResponsibility}).
  \item[D-REVISION-LIVE] The three underlying claim cards
    (\texttt{GLOSA-CC-20260904-0201/0202/0203}) are at revision 2 of an
    actively edited project; \texttt{status: Draft} and
    \texttt{independent\_check.status: NONE} on all three reflect that the
    founder has not yet signed off on any specific card instance, even though
    the founder has already ruled on genre (BBL-097) and hypothesis selection
    (BBL-095) at the project level.
  \item[D-NO-EPISTEMIC-VETO] No agency --- including any AI vendor, reviewer,
    or institution named in this paper --- holds an epistemic veto over H1,
    H2, or H3 merely by title, credential, or infrastructure ownership; any
    agency may propose, use, test, challenge, fork, translate, revise,
    restrict, or withdraw this work without needing permission, provided
    lineage is preserved (per the repo's own knowledge-validation stance,
    \texttt{CLAUDE.md}/\texttt{AGENTS.md}).
  \item[D-CANDIDATE-STATUS] H1, H2, and H3 are candidate hypotheses drafted by
    AI; the act of selecting among them (all three, per BBL-2026-09-04-095)
    was the founder's alone (\Cref{sec:standpoint-and-ownership}). No claim
    card here asserts candidate status was itself a form of testing.
  \item[D-EXTERNAL-INPUT] The lens this paper's readout-translation and
    hypothesis-drafting steps were performed under is named, cited, and
    linked throughout (\Cref{sec:lens}); naming it credits its source and
    lets a reader replace it --- it does not make the lens, or any hypothesis
    drafted under it, true.
\end{disclaimers}


\subsection{Knowledge state}
\label{sec:kstate}
This paper's overall knowledge state is \textbf{K0} (internal working note): every load-bearing
claim in it, including \Cref{sec:demonstration}'s findings, shares this K0 state. It is not K1
(public-provisional with a named independence class per claim and no class-5 human check yet)
because the three underlying claim cards' own \texttt{status} field is still \texttt{Draft} and
\texttt{independent\_check.status: PENDING} (checked by three cross-vendor routes and revised, not
yet re-checked by a further route or founder-approved, \Cref{sec:self-review}) --- the founder has
ruled on genre and selection at the project level but has not yet signed off on any specific card
instance (\Cref{sec:claim-matrix}).

\section{AI-assistance disclosure}
\label{sec:ai-disclosure}
This paper is a human--AI co-production, not AI-only automation and not human-only authorship.
A generative-AI assistant (vendor recorded in the local cooking log, not named here per gate rule 9) assisted with literature retrieval
and card drafting, the three claim cards' field content, the conversation-with-the-literature
tables, and this \LaTeX{} build. \textbf{AI output is not treated as evidence.} Every claim above
is accepted only when it matches a retained source record (a citation card's
\texttt{exact\_passage}, the Blackbox Log's own entries, a mechanically executed command's own
output) or a machine-checkable result reported in \Cref{sec:demonstration}. The human holds the
standpoint, the falsifier judgment, verification of load-bearing citations at the founder's own
discretion, and the final public commitment (\textbf{AIContribution $\neq$
EpistemicResponsibility}).

% ===========================================================================
\section{What would change our mind}
\label{sec:change-mind}
Per hypothesis, stated plainly rather than left implicit:

\begin{itemize}
\item \textbf{H1} would be abandoned if a genuine I5 reviewer, working from a schema-passing claim
card, found a false answer at Q2 or Q3 --- showing the mechanical check gives false reassurance
rather than tracking the property it names. It would be strengthened, not proven, by a corpus
comparison at I3+ (with a subset at I5) against a matched non-card corpus.
\item \textbf{H2} would be abandoned if a reader study found that credentialed audiences (editors,
grant reviewers, a practitioner's own community) still required an institutional signal even after
being shown a complete, well-formed claim card --- i.e.\ the credential doing load-bearing work the
card's fields could not substitute for. The Scholarly Kitchen credentialism case
(\texttt{cite-rigour-without-infrastructure-004}) is already the closest real instance of exactly
this challenge and is not explained away.
\item \textbf{H3} is untested; the one source that actually ran its comparison
(\texttt{cite-crossmodel-codereview-002}) found a direction-dependent, asymmetric effect rather than
a uniform cross-vendor-beats-same-vendor pattern. A seeded-error benchmark that reproduces that
asymmetry at the claim-card Q3/Q4 grain would be evidence against H3 as a general, direction-agnostic
claim; a benchmark that finds a uniform, direction-independent rise would be evidence for it.
\end{itemize}

\section{Conclusion}
\label{sec:conclusion}
This paper states three falsifiable propositions about whether a claim-card discipline can
substitute for institutional infrastructure in making a standalone scholar's claim reader-checkable,
places each against verified literature rather than memory or silence, and then runs the same
discipline on its own construction --- finding, honestly, a validator warning, a caught set of
composite quotes that changed the method's own citation rule mid-project, a router/manifest
mismatch the founder had to catch by hand, and a small stance-counting discrepancy in the
project's own summary statistics. None of this is smoothed into a clean success story: every claim
here stays \DrTier{}, every underlying card stays K0 and unsigned, and H3 stays explicitly
untested. What is offered is not a demonstrated result but a specified mechanism, an honest account
of what happened when it was pointed at itself, and a named falsifier for each of the three ways it
could be shown wrong.

\appendix

% ===========================================================================
\section{Claim matrix}
\label{sec:claim-matrix}
Per Chair Ruling v1 \S B2, this is a print-space compression of the three full claim cards
(\texttt{projects/GLS-2026-001\_rigour-without-infrastructure/claims/GLOSA-CC-20260904-0201/0202/0203.yaml});
field-presence here is mechanically checkable, correctness is what the independent check (Q5) exists
for.

\begin{table*}[t]
\centering
\scriptsize
\begin{tabularx}{\linewidth}{@{}p{0.035\linewidth} X p{0.14\linewidth} p{0.05\linewidth} p{0.08\linewidth} p{0.07\linewidth}@{}}
\toprule
\textbf{ID} & \textbf{Informal statement} & \textbf{Evidence / support} & \textbf{Tier} & \textbf{Produced by} & \textbf{Responsible} \\
\midrule
C1 & H1: if every claim is recorded as a claim card tied to E-A-D and passed through the independence ladder before release, then a standalone scholar with no institutional affiliation \emph{could, as a design proposition,} co-produce checkable-rigour knowledge with AI. & 15 verified citation cards (\Cref{tab:h1-lit}); \texttt{GLOSA-CC-20260904-0201} & \DrTier{} & joint (data$\to$inference); human (inference$\to$claim) & human \\
C2 & H2: rigour is proposed here, as a conceptual contrast, to be \emph{at least in part} a property a single claim can carry in its own recorded structure, not conferred solely from outside by institutional machinery. & 16 verified citation cards (\Cref{tab:h2-lit}); \texttt{GLOSA-CC-20260904-0202} & \DrTier{} & joint (data$\to$inference); human (inference$\to$claim) & human \\
C3 & H3: cross-vendor (I3) checking of undisclosed AI-fill/assumptions \emph{would rise} measurably compared to same-vendor (I0/I1) re-checking. \textbf{Untested; no seeded-error benchmark has been run by this project (n=0).} & 17 verified citation cards (\Cref{tab:h3-lit}), all analogical, none a direct test of Q3/Q4 detection; \texttt{GLOSA-CC-20260904-0203} & \DrTier{} & joint (data$\to$inference); human (inference$\to$claim) & human \\
\bottomrule
\end{tabularx}
\caption{Claim Matrix. All three cards: \texttt{k\_state: K0}, \texttt{status: Draft},
\texttt{independent\_check.status: PENDING} (checked by cross-vendor routes R1--R3 and revised,
\Cref{sec:self-review}; not yet re-checked by a further route or founder-approved),
\texttt{independence\_class: I3} (the citation-verification route, not the card's own sign-off).}
\label{tab:claim-matrix}
\end{table*}

\subsection{Five-questions detail}
\begin{table*}[t]
\centering
\scriptsize
\begin{tabularx}{\linewidth}{@{}p{0.035\linewidth} X X X p{0.1\linewidth}@{}}
\toprule
\textbf{ID} & \textbf{seen} & \textbf{ai\_filled} & \textbf{assumed} & \textbf{tested (indep.)} \\
\midrule
C1 & Blackbox Note L1; founder's abstract framing & candidate hypothesis text, evidence retrieval, card drafting (all disclosed, \texttt{disclaimers\_emitted: D-AIFILL}) & no source tests the full combined H1 mechanism as one system & I3 (citation checks); card itself unsigned \\
C2 & same L1; H2's own narrower conceptual framing & same AI-fill fields as C1 & the two SUPPORTS sources are same-lineage; cannot alone move H2 past \DrTier{} & I3; card itself unsigned \\
C3 & same L1; H3's cross-vendor detection-rate framing & same AI-fill fields as C1 & no seeded benchmark has been run; the closest analogous study found a direction-dependent effect & I3 (analogical citations only); card itself unsigned \\
\bottomrule
\end{tabularx}
\caption{Five-questions detail keyed to \Cref{tab:claim-matrix}. Full cards:
\texttt{projects/GLS-2026-001\_rigour-without-infrastructure/claims/GLOSA-CC-20260904-0201.yaml}
(and \texttt{-0202}, \texttt{-0203}).}
\label{tab:five-questions}
\end{table*}

% ===========================================================================
\section{Blackbox Note: how this work was made}
\label{sec:blackbox}
\emph{(Thai project term: ``Bantuek Klong Dam''.)} Curated verbatim lines --- typos kept, never
polished or translated in place --- followed by the cooking log. Full record archived separately
(\texttt{projects/GLS-2026-001\_rigour-without-infrastructure/blackbox\_note.yaml}, note id
\texttt{BB-2026-09-04-01}; public curated source
\texttt{blackbox/BLACKBOX\_NOTE\_glosa-paper\_2026-09-04.md}; public log entries cited by id below,
\texttt{blackbox/log/entries.jsonl}, concept DOI \texttt{10.5281/zenodo.22302518}).

\begin{table*}[t]
\centering
\scriptsize
\begin{tabularx}{\linewidth}{@{}p{0.09\linewidth} p{0.08\linewidth} X p{0.22\linewidth}@{}}
\toprule
\textbf{ID} & \textbf{Kind} & \textbf{English gloss (AI assistant's translation, tier \DrTier{}; effect on this paper)} & \textbf{Pointer} \\
\midrule
BBL-083 & proposal & ``the word `observation' is abstract, it is a verb; if we started from the problem instead, is there a process like this? We propose focusing on the problem first.'' --- separates ``observation'' (abstract verb) from ``problem state'' (concrete pre-existing thing); became FOUNDATION \S2.1a (\Cref{sec:problem}) & DOI \texttt{10.5281/zenodo.22302518}, entry \texttt{BBL-2026-09-04-083} \\
BBL-086/088 & finding / correction & ``if the problem is still ours, the question is still ours, and the selection of the hypothesis is still ours, and it leads to solving our problem, then use AI, use a cat, use a bear, go ahead --- because that is the goal of knowledge: to solve a problem for someone.'' (088 corrects 086) --- fixes which spine nodes must stay human: problem, question, and the \emph{selection} of the hypothesis; became FOUNDATION \S2.1c & DOI \texttt{10.5281/zenodo.22302518}, entry \texttt{BBL-2026-09-04-088} (corrects 086) \\
BBL-089 & ruling & Ratifies the finer-grained spine diagram wording (AI-drafted, human-ratified) --- became the spine diagram, FOUNDATION \S2.1a & \texttt{blackbox/log/entries.jsonl}, entry \texttt{BBL-2026-09-04-089} \\
BBL-091 & instruction & Install and use the skill to build this article in full --- became task B1/B2 launch & \texttt{blackbox/log/entries.jsonl}, entry \texttt{BBL-2026-09-04-091} \\
BBL-095 & ruling & ``bring all three together'' --- selects all three hypotheses (H1, H2, H3); became \texttt{hypothesis\_selection.yaml} & \texttt{blackbox/log/entries.jsonl}, entry \texttt{BBL-2026-09-04-095} \\
BBL-097 & ruling & ``make this document an academic article, a concept paper'' --- sets genre to CONCEPTUAL, overriding the router's \texttt{systematic\_review} output; became \texttt{genre\_decision.yaml} (\Cref{sec:genre}) & \texttt{blackbox/log/entries.jsonl}, entry \texttt{BBL-2026-09-04-097} \\
BBL-098 & instruction & Requires colliding own prior papers against world literature, not silence --- became own-lineage rows, \Cref{sec:literature} & \texttt{blackbox/log/entries.jsonl}, entry \texttt{BBL-2026-09-04-098} \\
BBL-099 & ruling & English only, no Thai edition; arXiv two-column theme, polished --- became this document's language and theme & \texttt{blackbox/log/entries.jsonl}, entry \texttt{BBL-2026-09-04-099} \\
BBL-100 & ruling & ``should we make it a hard rule: if a source is used it must carry a page and line locator plus a full link, never written from memory --- that reduces the problem at the source, before the check.'' --- mandates locator+link, source-first citation, after I3 caught composite quotes; became rule 17, \S7.8 (\Cref{sec:demonstration}) & DOI \texttt{10.5281/zenodo.22302518}, entry \texttt{BBL-2026-09-04-100} \\
BBL-101 & instruction & ``update it into the system'' --- adopts the citation rule immediately; became rule 17 adoption & \texttt{blackbox/log/entries.jsonl}, entry \texttt{BBL-2026-09-04-101} \\
\bottomrule
\end{tabularx}
\caption{Blackbox Note --- every row is id / English gloss / pointer. Verbatim Thai lines are kept
only in the cited records (the Blackbox Log, DOI \texttt{10.5281/zenodo.22302518}, or
\texttt{blackbox/log/entries.jsonl}), by founder ruling \texttt{BBL-2026-09-04-103} (English-only
paper); the glosses shown here are the AI assistant's translations, not the founder's own words,
tier \DrTier{}.}
\label{tab:blackbox-lines}
\end{table*}

\begin{table*}[t]
\centering
\scriptsize
\begin{tabularx}{\linewidth}{@{}p{0.1\linewidth} p{0.09\linewidth} p{0.12\linewidth} X X@{}}
\toprule
\textbf{Step} & \textbf{By} & \textbf{Input lines} & \textbf{Output} & \textbf{What changed} \\
\midrule
lens\_in & ai-assistant & L1 & \texttt{lens\_translation.md} & Drafted \texttt{question\_readout}, \texttt{local\_contrast\_space\_X}, \texttt{access\_relation\_R}, \texttt{claim\_function\_Phi\_z0} from L1 under the Readout Universe lens; not yet independently reviewed. \\
lens\_out & ai-assistant & L1--L13 & \texttt{hypotheses.md} & Drafted H1/H2/H3 as candidate world-language hypotheses; selection reserved to and performed by the founder (BBL-095). \\
lit\_review & ai-assistant & H1/H2/H3 & 48 citation cards, 3 \texttt{litreview\_manifest.yaml} & Ran targeted/scoping searches per hypothesis (not systematic), froze manifests, computed diversity/accuracy audits. \\
claim\_card & ai-assistant & hypotheses + literature & \texttt{GLOSA-CC-20260904-0201/0202/0203} & Filled full-shape cards; validated \texttt{PASS\_WITH\_LIMITS}; not yet founder-signed. \\
paper\_build & ai-assistant & all of the above & this document & Drafted the concept paper's prose, tables, and \LaTeX{} build; disclosed throughout, never authored. \\
\bottomrule
\end{tabularx}
\caption{Cooking log (Thai project term: ``karn pung'') --- append-only.}
\label{tab:cooking-log}
\end{table*}

\section{Human Mastery Gate}
\label{sec:mastery-gate}
Before submission the author confirms they can defend, without AI assistance: (1) the problem
(\Cref{sec:problem}) (2) the major literature families (\Cref{sec:literature}) (3) the closest
constructs (own-lineage rows, \Cref{sec:literature}) (4) the precise inadequacy (the contaminated
concept named in \Cref{sec:problem}) (5) the proposed mechanism (H1, \Cref{sec:propositions}) (6)
the boundary conditions ($E \wedge A \wedge D$, \Cref{sec:lens}) (7) the strongest counterargument
(the Scholarly Kitchen credentialism case, \Cref{tab:h2-lit}) (8) the strongest counterexample (the
direction-dependent cross-model result against H3, \Cref{tab:h3-lit}) (9) the falsifying record for
each proposition (\Cref{sec:propositions}) (10) the discriminating empirical test (H3's
not-yet-run seeded-error benchmark). \textbf{Gate status: PASS WITH NAMED GAPS} --- the founder has
not yet performed the class-5 (I5) sign-off on any of the three underlying claim cards
(\Cref{sec:claim-matrix}); this gap is named, not silently dropped.

\section{Provenance and lineage}
\label{sec:provenance}
Status vocabulary: \textsc{PRESERVE\_EXACT}, \textsc{PRESERVE\_FUNCTION}, \textsc{EXPAND},
\textsc{SUPERSEDE}. This paper \textsc{PRESERVE\_FUNCTION}s the claim-card schema and independence
ladder from \texttt{design/FOUNDATION\_v0.5.md} and \textsc{EXPAND}s it with a three-proposition
literature conversation and a self-application demonstration specific to this project
(GLS-2026-001). Direct ancestor: the Standalone Scholar dual-track architecture
\citep{cite-own-standalone-scholar-h1-005}. Core epistemic engine this paper's disclosure apparatus
is built on: the Readout Condition \citep{cite-own-readout-condition-h1-002}.

\balance
% References set \sloppy + \raggedright so long DOI/URL entries break at the
% column edge instead of producing overfull \hbox warnings (natbib/plainnat
% does not hyphenate URLs on its own); url package above loads with
% breaklinks so hyperref-wrapped URLs are also allowed to break mid-string.
\sloppy
\raggedright
\bibliographystyle{plainnat}
\bibliography{refs}

\end{document}
